% ===========================================================================
%  DAIM Research Report Series -- report template
%
%  1. Fill in the metadata block below.
%  2. Replace the placeholder text in each section.
%  3. Build with:  latexmk -pdf main.tex
%
%  Add [anonymous] to hide the author block for blind review, or [nologo] to
%  drop the masthead logo:  \documentclass[anonymous]{daimreport}
% ===========================================================================
\documentclass{daimreport}

% ---------------------------------------------------------------------------
%  Metadata -- everything the title page, running head and imprint need
% ---------------------------------------------------------------------------
\reporttitle{Report Title in Sentence Case}
\reportsubtitle{An optional subtitle giving the scope or the claim}
\shorttitle{Short Title for the Running Head}

% Use \daimmark{...} for the superscript markers, and explain them in
% \reportauthornote. Separate authors with commas.
\reportauthor{First A. Author\daimmark{*}, Second B. Author,
              Third C. Author\daimmark{\dag}}

\reportaffiliation{Centre of Excellence for Data Science, Artificial Intelligence
                   and Modelling (DAIM), University of Hull, United Kingdom}

\reportauthornote{\daimmark{*}Corresponding author.
                  \daimmark{\dag}Also at Department Name, Institution Name.}

\reportemail{first.author@hull.ac.uk}
\reportnumber{DAIM-RR-2026-000}

% Optional overrides -- leave commented out to accept the defaults.
%
% The masthead date is the compile date automatically, so it tracks whoever
% rebuilds the file. Uncomment \reportdate to pin it once the report is issued.
% \reportdate{16 March 2026}
% \reportseries{DAIM Research Report Series}
% \reportrights{\textcopyright\ 2026 University of Hull. All rights reserved.}
% \reportlogowidth{50mm}
\reportlogo{assets/DAIM_UoH_Logo.png}

\begin{document}

\makedaimtitle

% ---------------------------------------------------------------------------
%  Abstract -- one paragraph, no heading, set bold across the full measure.
%  Aim for 150-250 words: what the problem is, what you did, what you found.
% ---------------------------------------------------------------------------
\begin{daimabstract}
State the problem in one or two sentences, then say what this report
contributes. Summarise the approach without method detail, give the headline
result with a number if you have one, and close with what follows from it. The
abstract carries no heading and is set in bold across the full measure, so keep
it to a single paragraph. Readers who go no further should still come away
knowing what was done and what it means.
\end{daimabstract}

% Optional; delete if the series does not use keywords.
\daimkeywords{first keyword; second keyword; third keyword; fourth keyword}

% ---------------------------------------------------------------------------
%  Contents -- OFF by default, which is the convention this layout follows: a
%  published paper carries no printed contents list and navigates by the PDF
%  bookmarks, which are generated either way.
%
%  Uncomment the line below only for a long report (roughly fifteen pages up).
%  It is set the way a paper would set it -- no dotted leaders, tight leading,
%  page numbers flush right -- not the way a thesis would.
% ---------------------------------------------------------------------------
% \daimcontents

\section{Introduction}

Open with the problem and why it matters, in prose rather than bullets. Cite as
you go: parenthetically with \citep{placeholder-article}, or in the running text
as \citet{placeholder-book} showed. Multiple sources go in one call,
\citep{placeholder-article,placeholder-inproceedings}, so that they collapse
into a single parenthesis.

Paragraphs are indented and carry no extra vertical space between them, which is
the convention this layout follows. Use a footnote for an aside that would
otherwise interrupt the argument.\footnote{Footnotes are set small below a short
rule. Keep them rare; anything load-bearing belongs in the body or an appendix.}

State the contribution explicitly. A short list is appropriate here even though
the surrounding text is prose:

\begin{enumerate}
  \item \textbf{First contribution.} One sentence saying what is new and what it
        enables.
  \item \textbf{Second contribution.} Another sentence in the same shape.
  \item \textbf{Third contribution.} And a third, if there is one.
\end{enumerate}

Use a key box sparingly --- once or twice in a report --- for the claim you most
want a skimming reader to take away.

\begin{daimkeybox}{Headline finding}
State the single most important result or position here, in two or three
sentences. This panel is deliberately conspicuous, so it loses its force if it
appears on every page.
\end{daimkeybox}

\subsection{Scope}

Say what the report does and does not cover, and how it relates to adjacent
work. Second-level headings are numbered and unpunctuated.

\subsubsection{A third level, if the material needs it}

Three levels are available. If you find yourself wanting a fourth, use a run-in
paragraph heading instead.

\paragraph{Run-in heading} A bold run-in heading followed by a full stop, then
the text on the same line. Good for naming a method, an assumption or a case
without spending a numbered heading on it.

\section{Method}

Describe the approach in enough detail to be reproducible. Display mathematics
is numbered and referenced by number, so that the text can point at it:
\begin{equation}
  y = f(x; \theta) + \varepsilon,
  \qquad \varepsilon \sim \mathcal{N}(0, \sigma^2),
  \label{eq:model}
\end{equation}
where $x$ is the input, $\theta$ the parameters and $\varepsilon$ the residual.
Refer back with Equation~\eqref{eq:model} rather than repeating it.

Figures are centred, with the caption below and labelled for cross-reference.
Replace the placeholder box below with \verb|\includegraphics| or a TikZ
picture; Figure~\ref{fig:placeholder} shows where it goes.

\begin{figure}[t]
\centering
% Replace this placeholder with your artwork, e.g.
%   \includegraphics[width=0.8\textwidth]{assets/your-figure.pdf}
\fcolorbox{DAIMRule}{DAIMAccentTint}{%
  \parbox[c][36mm][c]{0.78\textwidth}{\centering\color{DAIMMuted}\small
    Figure artwork goes here\\[3pt]
    \texttt{\textbackslash includegraphics[width=0.8\textbackslash textwidth]\{file\}}}}
\caption{A one- or two-sentence caption that says what the reader should see,
not merely what the axes are. The label and text are separated by a vertical
bar, and the whole caption is centred.}
\label{fig:placeholder}
\end{figure}

\section{Results}

Tables use \texttt{booktabs} rules with a tinted header row. Keep them free of
vertical rules, and put the caption below the table.

\begin{table}[t]
\centering
\begin{tabularx}{\textwidth}{L{28mm} R{22mm} R{22mm} Y}
\toprule
\daimheadrow \daimth{Condition} & \daimth{Metric A} & \daimth{Metric B} & \daimth{Notes} \\
\midrule
Baseline        & 0.00 & 0.00 & Reference condition. \\
Variant one     & 0.00 & 0.00 & One clause on what changed. \\
Variant two     & 0.00 & 0.00 & One clause on what changed. \\
\textbf{Ours}   & \textbf{0.00} & \textbf{0.00} & Bold the row you want read first. \\
\bottomrule
\end{tabularx}
\caption{Caption below the table, saying what the comparison establishes.
Define every abbreviation used in the header.}
\label{tab:placeholder}
\end{table}

Refer to floats by number --- Table~\ref{tab:placeholder} and
Figure~\ref{fig:placeholder} --- never as ``the table above'', since floats
move.

Use the neutral grey panel for verbatim material: a transcript, a prompt, a
worked example, a piece of a specification.

% Short panels read better kept whole; drop [breakable=false] for long ones.
\begin{daimbox}[breakable=false]{Example transcript}
\textbf{Input:} The prompt, query or input case, quoted verbatim.

\textbf{Output:} The response or result, quoted verbatim and truncated with an
ellipsis if it runs long.

\textit{A short italic gloss on what the example demonstrates.}
\end{daimbox}

\section{Discussion}

Interpret the results: what they support, what they do not, and where the
approach breaks down. Be specific about limitations --- a reader who finds them
for you will trust the rest less. Reports and preprints cite the same way as
anything else \citep{placeholder-report}.

\section{Conclusion}

Restate the contribution in a sentence, then say what should happen next. Avoid
introducing new claims here.

% ---------------------------------------------------------------------------
%  Back matter
% ---------------------------------------------------------------------------
\daimsection{Acknowledgements}

Funding, facilities and people who helped but are not authors. Name the grant
and its number if there is one. Use \verb|\daimsection| for headings that
should appear in the contents but carry no number.

\daimbibliography{references}

% ---------------------------------------------------------------------------
%  Appendices -- sections become "Appendix A.", floats renumber to A1, A2, ...
%  Delete everything below if the report has no appendices.
% ---------------------------------------------------------------------------
\daimappendix

\section{Notation}

Material that supports the argument but would interrupt it: notation,
derivations, full parameter tables, extra results.

\begin{table}[h]
\centering
\begin{tabularx}{0.75\textwidth}{L{22mm} Y}
\toprule
\daimheadrow \daimth{Symbol} & \daimth{Meaning} \\
\midrule
$x$        & Input variable \\
$y$        & Observed response \\
$\theta$   & Model parameters \\
$\varepsilon$ & Residual term \\
\bottomrule
\end{tabularx}
\caption{Notation used throughout the report.}
\label{tab:notation}
\end{table}

\section{Additional results}

A second appendix, if you need one.

\end{document}
